\documentclass[a4paper,12pt]{article}
\usepackage[T2A]{fontenc}
\usepackage[utf8]{inputenc}
\usepackage[russian]{babel}
\usepackage{amsmath,amssymb,mathtools}
\usepackage{helvet}
\renewcommand{\sfdefault}{cmss}
\renewcommand{\familydefault}{\sfdefault}
\usepackage{icomma}
\usepackage[a4paper,left=23mm,right=23mm,top=23mm,bottom=23mm,headheight=15pt]{geometry}
\usepackage{microtype}
\usepackage{tikz}
\usetikzlibrary{calc,arrows.meta,positioning,shapes.geometric}
\usepackage{tabularx,booktabs}
\usepackage{enumitem}
\usepackage{fancyhdr}
\usepackage{setspace}
\usepackage{xcolor}
\usepackage{hyperref}
\hypersetup{hidelinks,pdftitle={Чек-лист. Логарифмы и преобразования},pdfauthor={Лёвин Артём Александрович}}
\definecolor{accent}{HTML}{137A7A}
\definecolor{darkaccent}{HTML}{0A525A}
\definecolor{textgray}{HTML}{2F383B}
\definecolor{rulegray}{HTML}{C5D4D4}
\definecolor{pale}{HTML}{EEF4F3}
\color{textgray}
\onehalfspacing
\setlength{\parindent}{0pt}
\setlength{\parskip}{4pt}
\setcounter{secnumdepth}{0}
\setlist[enumerate]{leftmargin=1.9em,itemsep=5pt,topsep=5pt}
\pagestyle{fancy}
\fancyhf{}
\fancyhead[L]{\small\color{darkaccent}Марк Гукин\quad\textbar\quad 18.10.2026}
\fancyhead[R]{\small\color{textgray}Логарифмы}
\fancyfoot[C]{\small\color{textgray}\thepage}
\renewcommand{\headrulewidth}{0.35pt}
\renewcommand{\footrulewidth}{0pt}
\newcommand{\sect}[1]{\par\vspace{4mm}{\Large\bfseries\color{darkaccent}#1}\par\vspace{1.5mm}\color{textgray}}
\newcommand{\subj}[1]{\par\vspace{2mm}{\large\bfseries\color{darkaccent}#1}\par\color{textgray}}
\newcommand{\key}[1]{\textbf{\color{darkaccent}#1}}
\newcommand{\checkitem}[1]{\par\noindent$\square$\enspace #1\par}
\begin{document}

% Титульный лист.
\thispagestyle{empty}
\begin{center}
\vspace*{36mm}
{\small\bfseries\color{darkaccent}ИНДИВИДУАЛЬНОЕ УЧЕБНОЕ ПОСОБИЕ}\par
\vspace{19mm}
{\fontsize{34}{40}\selectfont\bfseries\color{darkaccent}Чек-лист}\par
\vspace{9mm}
{\Large\bfseries Логарифмы}\par
\vspace{2mm}
{\large Свойства и преобразования выражений}\par
\vspace{22mm}
{\large\bfseries Марк Гукин}\par
\vspace{8mm}
{\normalsize 18 октября 2026 года}\par
\vfill
{\normalsize Лёвин Артём Александрович}\par
{\normalsize\bfseries\color{darkaccent}эксклюзивно для Марка Гукина}\par
\vspace{18mm}
\begin{tikzpicture}[x=1cm,y=1cm]
  \draw[darkaccent,opacity=.45,line width=.9pt]
  (-7.4,0) .. controls (-4.9,.28) and (-3.5,-.38) .. (-.3,.13)
  .. controls (2.4,.50) and (4.4,-.32) .. (7.4,.05);
\end{tikzpicture}
\vspace{8mm}
\end{center}
\clearpage

\sect{01\quad Что означает логарифм}
\[\boxed{\log_a b=x\quad\Longleftrightarrow\quad a^x=b}\]
Читаем: \emph{логарифм числа $b$ по основанию $a$} --- это показатель степени, в которую нужно возвести $a$, чтобы получить $b$.
\[
\boxed{a>0,\quad a\ne1,\quad b>0.}
\]
В записи $\log_a b=x$: $a$ --- \key{основание}, $b$ --- \key{аргумент} (тело логарифма), $x$ --- \key{значение}.
\[
\log_2 16=4\ (2^4=16),\qquad
\log_a a=1,\qquad \log_a1=0.
\]

\sect{02\quad Показатели степеней и логарифмы}
\subj{Степень у аргумента}
\[
\boxed{\log_a(b^n)=n\log_a b.}
\]
Показатель можно вынести \emph{без изменения}. Формула работает и в обратном направлении: $n\log_a b=\log_a(b^n)$.
\[
\log_2 16=\log_2 2^4=4\log_2 2=\boxed{4}.
\]
\[
\log_2\frac{1}{\sqrt[5]{16}}
=\log_2(2^{-4/5})
=-\frac45\log_2 2=\boxed{-\frac45}.
\]
\subj{Степень у основания}
\[
\boxed{\log_{a^m}(b^n)=\frac nm\log_a b,\qquad m\ne0.}
\]
\key{Правило:} показатель \emph{основания} выходит обратным множителем, показатель \emph{аргумента} --- прямым множителем.
\[
\log_{1/\sqrt2}8=\log_{2^{-1/2}}2^3
=\frac{3}{-1/2}=\boxed{-6}.
\]
\[
\log_{1/9}\frac1{\sqrt[7]{81}}
=\log_{3^{-2}}3^{-4/7}
=\frac{-4/7}{-2}=\boxed{\frac27}.
\]
Все равенства на этой странице используют положительные основания и аргументы. В общей формуле также $a\ne1$.

\clearpage
\sect{03\quad Сумма, разность, множитель}
Для одного основания $a$ и положительных $b,c$:
\[
\boxed{\log_a b+\log_a c=\log_a(bc)}
\]
\[
\boxed{\log_a b-\log_a c=\log_a\frac bc}
\]
Перед сворачиванием логарифмов нужно сделать \key{основания одинаковыми}, а числовые коэффициенты перенести в степень аргумента.
\[
\begin{aligned}
2\log_2 3+\log_2 5
&=\log_2(3^2)+\log_2 5\\
&=\log_2(9\cdot5)=\boxed{\log_2 45}.
\end{aligned}
\]
\[
\log_2 3+\log_2 5-\log_2 7
=\boxed{\log_2\frac{15}{7}}.
\]
Если основания различаются, сначала приведите их к одному:
\[
\begin{aligned}
\log_2 3+\log_8 5
&=\log_2 3+\frac13\log_2 5\\
&=\boxed{\log_2\bigl(3\sqrt[3]{5}\bigr)}.
\end{aligned}
\]
\key{Не перепутайте:} произведение аргументов соответствует \emph{сумме} логарифмов; отношение аргументов соответствует \emph{разности}.

\sect{04\quad Основное логарифмическое тождество}
\[
\boxed{a^{\log_a b}=b}
\]
Основание степени должно совпасть с основанием логарифма. Если оно не совпадает, сначала преобразуйте степень или логарифм.
\[
2^{\log_8 3}
=2^{\frac13\log_2 3}
=\boxed{\sqrt[3]{3}}.
\]
\[
8^{\log_2 3}
=(2^3)^{\log_2 3}
=(2^{\log_2 3})^3=\boxed{27}.
\]
\[
\begin{aligned}
2^{1+\log_2 5}&=2\cdot 2^{\log_2 5}=\boxed{10},\\
2^{1-\log_2 5}&=\frac{2}{2^{\log_2 5}}=\boxed{\frac25}.
\end{aligned}
\]
\[
3^{2+\log_{27}2}=3^2\cdot3^{\frac13\log_3 2}
=\boxed{9\sqrt[3]{2}}.
\]

\clearpage
\sect{05\quad Число в логарифмической записи}
Основное тождество можно читать справа налево. Для $a>0$, $a\ne1$ и любого действительного $t$:
\[
\boxed{t=\log_a(a^t).}
\]
Например, одно и то же число $5$ можно представить с разными основаниями:
\[
5=\log_{10}(10^5)=\log_2(2^5)=\log_\pi(\pi^5).
\]
А для любого $b>0$:\quad $b=a^{\log_a b}$.
\[
 e=\pi^{\log_\pi e},\qquad
 \pi=e^{\ln\pi}.
\]
Здесь $\ln t=\log_e t$.

\sect{06\quad Переход к новому основанию}
Для $a,c>0$, $a,c\ne1$ и $b>0$:
\[
\boxed{\log_a b=\frac{\log_c b}{\log_c a}.}
\]
Новое основание $c$ можно выбрать удобным для вычислений:
\[
\log_2 3=\frac{\log_{10}3}{\log_{10}2}
=\frac{\ln3}{\ln2}.
\]
\subj{Обращение логарифма}
При $b\ne1$ допустимо поменять основание и аргумент:
\[
\boxed{\log_a b=\frac1{\log_b a}.}
\]
\key{Пример из задач экзаменационного формата.} Известно, что $\log_b a=-2$. Найдите $\log_a(b^6\sqrt a)$.
\[
\begin{aligned}
\log_a(b^6\sqrt a)
&=6\log_a b+\frac12\log_a a\\
&=\frac6{\log_b a}+\frac12
=\frac6{-2}+\frac12=\boxed{-\frac52}.
\end{aligned}
\]
В этом примере $a,b>0$, $a\ne1$, $b\ne1$.

\clearpage
\sect{07\quad Дополнительное тождество}
Для $a,c>0$ и $b>0$, $b\ne1$:
\[
\boxed{a^{\log_b c}=c^{\log_b a}.}\tag{\(*\)}
\]
Эта формула полезна при перестановке основания степени и аргумента логарифма. Для развёрнутого решения задачи повышенной сложности её можно предварительно вывести.

\subj{Короткое доказательство}
Пусть $p=\log_b a$, $q=\log_b c$. Тогда $a=b^p$, $c=b^q$, следовательно,
\[
 a^{\log_b c}=(b^p)^q=b^{pq}=(b^q)^p=c^{\log_b a}.
\]
Таким образом, обе части равны $b^{pq}$.

\sect{08\quad Объединяем методы}
\key{Пример.} Вычислите
\[
\frac{\log_4 125}{\log_8 625}.
\]
Все основания и аргументы записываем степенями:
\[
4=2^2,\quad 8=2^3,\quad 125=5^3,\quad 625=5^4.
\]
Тогда
\[
\frac{\log_{2^2}5^3}{\log_{2^3}5^4}
=\frac{\frac32\log_2 5}{\frac43\log_2 5}
=\frac32\cdot\frac34=\boxed{\frac98}.
\]
Сокращение на $\log_2 5$ допустимо, потому что $\log_2 5\ne0$.

\subj{Проверочный список}
\checkitem{Определены основания и аргументы логарифмов; учтены $a>0$, $a\ne1$, $b>0$.}
\checkitem{Корни, обратные величины и составные числа представлены степенями.}
\checkitem{При выносе степени основания использован \emph{обратный} показатель.}
\checkitem{Перед сложением и вычитанием логарифмы приведены к одному основанию.}
\checkitem{При переходе к дроби проверено, что знаменатель не обращается в нуль.}
\checkitem{Степени, знаки и итоговые значения проверены подстановкой.}

\clearpage
\sect{Алгоритм: работа с логарифмами}
\vspace{9mm}
\begin{center}
\begin{tikzpicture}[
 node distance=10mm,
 terminal/.style={ellipse,draw=darkaccent,fill=pale,minimum width=38mm,minimum height=12mm,align=center,line width=.85pt},
 action/.style={rounded corners=2pt,draw=darkaccent,minimum width=120mm,text width=111mm,minimum height=17mm,align=center,inner sep=4.5pt,line width=.85pt},
 arrow/.style={-{Stealth[length=3mm]},line width=.85pt,draw=darkaccent}
]
\node[terminal] (s) {Начало};
\node[action,below=of s] (one) {Проверьте допустимость каждого логарифма: основание положительно и отлично от $1$, аргумент положителен};
\node[action,below=of one] (two) {Представьте числа, обратные величины и корни как степени удобных оснований};
\node[action,below=of two] (three) {Перенесите показатели: из аргумента --- напрямую, из основания --- обратным множителем};
\node[action,below=of three] (four) {Если логарифмов несколько, приведите их к одному основанию и перенесите коэффициенты в аргументы};
\node[action,below=of four] (five) {Сверните суммы и разности либо примените основное логарифмическое тождество};
\node[action,below=of five] (six) {Сократите ненулевые множители, вычислите результат и проверьте знаки};
\node[terminal,below=of six] (e) {Ответ};
\foreach \x/\y in {s/one,one/two,two/three,three/four,four/five,five/six,six/e}{\draw[arrow](\x)--(\y);}
\end{tikzpicture}
\end{center}
\clearpage

\sect{Тренировочный блок\quad\textnormal{Путь А}}
\textit{10 заданий. Решайте последовательно. Записывайте преобразования и конечный ответ.}
\vspace{3mm}
\begin{enumerate}[label=\textbf{\arabic*.},itemsep=12pt,parsep=0pt]
\item Вычислите: $\log_5 125$.
\item Вычислите: $\displaystyle\log_2\frac1{\sqrt[3]{16}}$.
\item Вычислите: $\displaystyle\log_{1/\sqrt3}27$.
\item Вычислите: $\displaystyle\log_{1/9}\frac1{\sqrt[5]{81}}$.
\item Вычислите: $\displaystyle\log_7 14+\log_7\frac12$.
\item Вычислите: $\displaystyle2\log_3 5+\log_3 9-\log_3 25$.
\item Вычислите: $\displaystyle2^{\log_8 125}$.
\item Вычислите: $\displaystyle3^{2+\log_3 4}$.
\item Известно, что $a,b>0$, $a\ne1$, $b\ne1$ и $\log_b a=-4$. Найдите
\[\log_a\bigl(b^6\sqrt a\bigr).\]
\item Вычислите:
\[\frac{\log_9 128}{\log_{27}256}.\]
\end{enumerate}
\vspace{5mm}
{\small\color{darkaccent}\textbf{Перед проверкой:}\quad
$\square$ основания\quad
$\square$ знаки\quad
$\square$ обратные показатели\quad
$\square$ ОДЗ.}
\clearpage

\sect{Ответы к тренировочному блоку}
\vspace{6mm}
\renewcommand{\arraystretch}{1.62}
\begin{tabularx}{\linewidth}{@{}cX@{}}
\toprule
\textbf{№} & \textbf{Ответ}\\
\midrule
1 & $3$\\
2 & $-\dfrac43$\\
3 & $-6$\\
4 & $\dfrac25$\\
5 & $1$\\
6 & $2$\\
7 & $5$\\
8 & $36$\\
9 & $-1$\\
10 & $\dfrac{21}{16}$\\
\bottomrule
\end{tabularx}
\vspace{13mm}
\subj{Короткая самопроверка}
\begin{enumerate}[label=\textbf{\arabic*.},itemsep=7pt]
\item В заданиях 2--4 приведите оба числа к одной основе и сравните показатели.
\item В заданиях 5--6 убедитесь, что коэффициенты перенесены в степени \emph{до} объединения логарифмов.
\item В задании 9 используйте $\log_a b=1/\log_b a$; в задании 10 сократите общий ненулевой множитель.
\end{enumerate}
\end{document}
